\chapter{Diseases of the Heart and Cardiovascular System}
\label{ch:heart}
The heart is a muscular organ that pumps blood through the circulatory system, delivering oxygen and nutrients to tissues and removing metabolic waste. Cardiovascular diseases remain the leading cause of death globally, responsible for approximately 18 million deaths annually. This chapter covers diseases of the heart muscle, valves, conduction system, coronary arteries, and peripheral vasculature.

%-------------------------------------------------------------
\section{Cardiac Arrhythmias}
%-------------------------------------------------------------

%-------------------------------------------------------------

%-------------------------------------------------------------

%-------------------------------------------------------------

%-------------------------------------------------------------

%-------------------------------------------------------------

%-------------------------------------------------------------

\label{sec:Cardiac Arrhythmias}
%-------------------------------------------------------------%-------------------------------------------------------------

\subsection{Atrial Fibrillation}
\label{sec:Atrial Fibrillation}

Atrial fibrillation is the most common sustained cardiac arrhythmia, affecting 2--4\% of the adult population and increasing in prevalence with age; risk factors include hypertension, heart failure, valvular disease, obesity, obstructive sleep apnea, hyperthyroidism, and alcohol excess. The condition results from disorganized electrical activity in the atria, leading to irregular and often rapid ventricular response. Clinical features include palpitations, fatigue, dyspnea, and decreased exercise tolerance; complications include thromboembolism (stroke), tachycardia-mediated cardiomyopathy, and impaired quality of life. Diagnosis is by ECG showing absent P waves, irregularly irregular R-R intervals, and variable QRS morphology. Management involves rate control (beta-blockers, calcium channel blockers, digoxin) and rhythm control (antiarrhythmic drugs, cardioversion, catheter ablation); the CHA$_2$DS$_2$-VASc score guides anticoagulation decisions, with direct oral anticoagulants preferred over warfarin for non-valvular atrial fibrillation. Prognosis depends on risk factor management and stroke prevention.

\subsection{Brugada Syndrome}
\label{sec:Brugada Syndrome}
Brugada syndrome is an inherited sodium channelopathy (SCN5A mutation most common) characterized by an ECG pattern of coved ST-segment elevation in the right precordial leads (V1--V2) and an increased risk of ventricular fibrillation and sudden cardiac death, particularly during sleep or at rest; it is more common in men and in Southeast Asian populations. The condition results from mutations in cardiac sodium channels leading to abnormal repolarization. Clinical features include syncope, nocturnal agitated resuscitation, or sudden cardiac death; patients may be asymptomatic. Diagnosis is based on ECG findings, with provocation testing using sodium channel blockers (ajmaline, flecainide) to unmask the pattern when necessary. Management involves implantable cardioverter-defibrillator placement for high-risk patients; fever and certain drugs should be avoided as they can exacerbate the ECG abnormality. Prognosis is excellent in low-risk patients and improved with implantable cardioverter-defibrillator in high-risk patients.

\subsection{Long QT Syndrome}
\label{sec:Long QT Syndrome}
Long QT syndrome is a group of inherited cardiac channelopathies characterized by prolongation of the QT interval on ECG and predisposition to torsades de pointes ventricular arrhythmias and sudden death; over 17 genes have been identified, with LQT1 (KCNQ1), LQT2 (KCNH2), and LQT3 (SCN5A) accounting for the majority. The condition results from mutations in cardiac ion channels affecting ventricular repolarization. Triggers vary by genotype: exercise (LQT1), auditory stimuli and emotional stress (LQT2), and rest or sleep (LQT3). Clinical features include syncope, seizures, and sudden cardiac death. Diagnosis is by ECG (corrected QT interval greater than 470 ms in men, greater than 480 ms in women is highly suggestive), genetic testing, and clinical criteria (Schwartz score). Management includes beta-blockers (nadolol preferred), avoidance of QT-prolonging drugs, and implantable cardioverter-defibrillator for high-risk patients; left cardiac sympathetic denervation is an option for drug-refractory cases. Prognosis is excellent with appropriate management and avoidance of triggers.

\subsection{Sick Sinus Syndrome}
\label{sec:Sick Sinus Syndrome}
Sick sinus syndrome (SSS) is a group of cardiac arrhythmias caused by dysfunction of the sinoatrial node, the heart's primary pacemaker, with an incidence increasing with age, affecting 1 in 600 patients over 65 years. The condition results from age-related degenerative fibrosis of the SA node, ischemia, infiltrative diseases (amyloidosis, sarcoidosis, hemochromatosis), cardiomyopathy, or surgical trauma. Iatrogenic causes include beta-blockers, calcium channel blockers, digoxin, and antiarrhythmic drugs. Clinical features include sinus bradycardia (often profound, <40 bpm), sinus arrest (pauses >3 seconds causing syncope), sinoatrial exit block, and the tachycardia-bradycardia syndrome (alternating paroxysmal atrial tachyarrhythmias, usually atrial fibrillation, with sinus bradycardia). Symptoms include fatigue, lightheadedness, presyncope, syncope (Stokes-Adams attacks), dyspnea on exertion, and palpitations. Diagnosis is based on 12-lead ECG showing bradyarrhythmia, 24-hour Holter monitoring (most useful for correlating symptoms with rhythm), event monitor, and for equivocal cases, electrophysiology study. Management involves permanent pacemaker implantation (dual-chamber or atrial pacing) for symptomatic SSS. Pacemaker therapy resolves symptoms and prevents syncope but does not reduce mortality. Rate-slowing medications for tachycardia-bradycardia syndrome may be used after pacemaker placement. Prognosis is good with appropriate pacing; SSS itself is not life-threatening but can cause significant morbidity from syncope and falls.

\subsection{Ventricular Tachycardia}
\label{sec:Ventricular Tachycardia}
Ventricular tachycardia is a potentially life-threatening arrhythmia originating from the ventricles at a rate of 100--250 beats per minute; it may be monomorphic or polymorphic. The most common cause is structural heart disease, particularly prior myocardial infarction with scar-related reentry. The condition results from abnormal electrical activity in ventricular myocardium. Hemodynamically unstable ventricular tachycardia requires immediate synchronized cardioversion. Clinical features include palpitations, dizziness, syncope, and cardiac arrest. Diagnosis is by ECG. Management of stable monomorphic ventricular tachycardia involves intravenous amiodarone or procainamide followed by catheter ablation; implantable cardioverter-defibrillators are indicated for secondary prevention (survivors of cardiac arrest) and primary prevention in high-risk patients; torsades de pointes is treated with intravenous magnesium. Prognosis depends on the underlying heart disease and effectiveness of prevention.

%-------------------------------------------------------------
\section{Cardiomyopathies}
%-------------------------------------------------------------

%-------------------------------------------------------------

%-------------------------------------------------------------

Cardiomyopathies are a heterogeneous group of myocardial disorders characterized by mechanical or electrical dysfunction of the heart muscle. These conditions frequently lead to inappropriate ventricular hypertrophy, dilation, or restrictive filling dynamics, impairing cardiac output. Etiologies encompass genetic mutations, ischemic injury, metabolic disorders, toxic exposure, and autoimmune inflammatory infiltration.

%-------------------------------------------------------------

%-------------------------------------------------------------

%-------------------------------------------------------------

%-------------------------------------------------------------

\label{sec:Cardiomyopathies}
%-------------------------------------------------------------%-------------------------------------------------------------

\subsection{Dilated Cardiomyopathy}
\label{sec:Dilated Cardiomyopathy}
Dilated cardiomyopathy is a condition characterized by ventricular chamber dilation and systolic dysfunction in the absence of coronary artery disease, valvular disease, or hypertensive heart disease. Causes include familial or genetic factors (mutations in titin, lamin A/C, and other cytoskeletal genes in 30--50\%), viral myocarditis, alcohol, toxins (doxorubicin), peripartum cardiomyopathy, and tachycardia-mediated cardiomyopathy. The condition results from impaired myocardial contractility leading to ventricular dilation. Clinical features include heart failure symptoms. Diagnosis is based on echocardiography showing a dilated left ventricle with reduced ejection fraction; endomyocardial biopsy may be indicated for suspected myocarditis. Management follows guideline-directed medical therapy for heart failure with reduced ejection fraction; anticoagulation may be considered for high thromboembolic risk; genetic counseling and family screening are recommended. Prognosis depends on the underlying cause and response to therapy.

\subsection{Hypertrophic Cardiomyopathy}
\label{sec:Hypertrophic Cardiomyopathy}
Hypertrophic cardiomyopathy is the most common inherited cardiovascular disease, affecting 1 in 500 individuals, characterized by asymmetric left ventricular hypertrophy (typically interventricular septum of 15 mm or greater) in the absence of other loading conditions; it is a leading cause of sudden cardiac death in young athletes. The condition results from mutations in sarcomere proteins (beta-myosin heavy chain, myosin-binding protein C being most common). The pathophysiology involves diastolic dysfunction, dynamic left ventricular outflow tract obstruction (present in 70\%), and myocardial ischemia. Clinical features include exertional dyspnea, chest pain, syncope, and palpitations. Diagnosis is by echocardiography and cardiac MRI. Management includes beta-blockers, calcium channel blockers, and mavacamten for obstructive disease; septal ablation or myectomy is for refractory obstruction; implantable cardioverter-defibrillator is placed for high-risk patients. Prognosis is generally good with appropriate management, though sudden cardiac death risk requires stratification.

\subsection{Myocarditis}
\label{sec:Myocarditis}
Myocarditis is an inflammation of the myocardium, most commonly caused by viral infection (Coxsackievirus B, adenovirus, parvovirus B19, SARS-CoV-2), but also by autoimmune diseases, drugs (immune checkpoint inhibitors), and toxins. The condition results from direct viral injury or immune-mediated myocardial damage. Clinical features range from mild chest pain and myalgias to fulminant heart failure, cardiogenic shock, and sudden death. Diagnosis involves elevated troponin, ECG changes (diffuse ST elevation), echocardiography showing regional or global wall motion abnormalities, and cardiac MRI with Lake Louise criteria; endomyocardial biopsy remains the gold standard for definitive diagnosis. Management is supportive, with guideline-directed medical therapy for heart failure; immunosuppression may be considered in autoimmune or giant cell myocarditis. Prognosis is generally favorable, with most patients recovering, but a subset develops dilated cardiomyopathy.

\subsection{Restrictive Cardiomyopathy}
\label{sec:Restrictive Cardiomyopathy}
Restrictive cardiomyopathy is a condition characterized by impaired ventricular filling due to increased myocardial stiffness with preserved or near-preserved systolic function; causes include amyloidosis (AL and ATTR types), sarcoidosis, hemochromatosis, endomyocardial fibrosis, and infiltrative diseases. The condition results from myocardial infiltration or fibrosis restricting ventricular filling. Clinical features include predominant right-sided heart failure (edema, ascites, hepatomegaly, jugular venous distension). Diagnosis is based on echocardiography showing small to normal-sized ventricles with biatrial enlargement and restrictive filling pattern; cardiac MRI demonstrates late gadolinium enhancement; endomyocardial biopsy confirms diagnosis. Management targets the underlying cause (tafamidis for ATTR amyloidosis, chemotherapy for AL amyloidosis) with diuretics for symptom management. Prognosis varies by etiology.

%-------------------------------------------------------------
\section{Coronary Artery Disease}
%-------------------------------------------------------------

%-------------------------------------------------------------

%-------------------------------------------------------------

%-------------------------------------------------------------

%-------------------------------------------------------------

%-------------------------------------------------------------

%-------------------------------------------------------------

\label{sec:Coronary Artery Disease}
%-------------------------------------------------------------%-------------------------------------------------------------

\subsection{Angina Pectoris}
\label{sec:Angina Pectoris}
Angina pectoris is chest discomfort caused by transient myocardial ischemia without myocardial necrosis. Stable angina occurs predictably with exertion or emotional stress and is relieved by rest or nitroglycerin, reflecting fixed coronary atherosclerotic stenosis; unstable angina occurs unpredictably or at rest and represents acute coronary syndrome without troponin elevation; microvascular angina results from coronary microvascular dysfunction. The condition results from an imbalance between myocardial oxygen supply and demand. Clinical features include chest pressure or discomfort that may radiate to the left arm or jaw. Diagnosis involves ECG (may show ST depression during episodes), stress testing, and coronary angiography for definitive anatomic assessment. Management of stable angina includes antianginals (beta-blockers, calcium channel blockers, nitrates), risk factor modification, and revascularization when medical therapy is inadequate; unstable angina requires urgent antiplatelet therapy, anticoagulation, and early invasive strategy. Prognosis depends on the extent of coronary disease and risk factor control.

\subsection{Ischemic Cardiomyopathy}
\label{sec:Ischemic Cardiomyopathy}
Ischemic cardiomyopathy is myocardial dysfunction resulting from chronic ischemia or prior myocardial infarction, leading to progressive ventricular dilation and heart failure; it is the most common cause of heart failure worldwide. The condition results from myocyte loss, fibrosis, and adverse ventricular remodeling following coronary artery disease. Clinical features include symptoms of heart failure: dyspnea on exertion, orthopnea, paroxysmal nocturnal dyspnea, fatigue, and peripheral edema. Diagnosis is based on echocardiography showing reduced left ventricular ejection fraction (less than 40\%), regional wall motion abnormalities, and mitral regurgitation. Management includes guideline-directed medical therapy (ACE inhibitors, ARBs, ARNI, beta-blockers, mineralocorticoid receptor antagonists, SGLT2 inhibitors, and diuretics); coronary revascularization may improve outcomes in selected patients; advanced therapies include cardiac resynchronization therapy, implantable cardioverter-defibrillators, and cardiac transplantation. Prognosis depends on the extent of myocardial dysfunction and response to therapy.

\subsection{Myocardial Infarction}
\label{sec:Myocardial Infarction}
Myocardial infarction is irreversible myocardial necrosis resulting from prolonged ischemia, most commonly caused by acute thrombotic occlusion of a coronary artery at the site of a ruptured or eroded atherosclerotic plaque. The condition results from complete or near-complete coronary artery occlusion. Clinical features include crushing chest pain, radiation to the left arm or jaw, diaphoresis, nausea, and dyspnea. Diagnosis is based on ECG (ST-elevation MI shows ST-segment elevation; non-ST-elevation MI presents without ST elevation but with elevated cardiac biomarkers) and elevated cardiac troponin. Management of STEMI requires emergent reperfusion therapy with primary percutaneous coronary intervention within 90 minutes or fibrinolysis within 30 minutes of first medical contact; treatment includes dual antiplatelet therapy, anticoagulation, beta-blockers, ACE inhibitors, and statins. Prognosis has improved with timely reperfusion and long-term secondary prevention.

%-------------------------------------------------------------
\section{Heart Failure}
%-------------------------------------------------------------

%-------------------------------------------------------------

%-------------------------------------------------------------

Heart failure is a complex clinical syndrome resulting from structural or functional cardiac impairment that compromises ventricular filling or ejection. Inadequate systemic perfusion activates compensatory neurohormonal pathways (renin-angiotensin-aldosterone system and sympathetic nervous system), driving progressive remodeling. Management combines lifestyle interventions, guideline-directed medical therapy, implantable devices, and advanced mechanical circulatory support.

%-------------------------------------------------------------

%-------------------------------------------------------------

%-------------------------------------------------------------

%-------------------------------------------------------------

\label{sec:Heart Failure}
%-------------------------------------------------------------%-------------------------------------------------------------

\subsection{Acute Decompensated Heart Failure}
\label{sec:Acute Decompensated Heart Failure}
Acute decompensated heart failure is a sudden worsening of heart failure signs and symptoms, most commonly presenting as pulmonary congestion with dyspnea, orthopnea, and crackles on auscultation; it is the leading cause of hospitalization in patients over 65. The condition results from precipitating factors including medication non-adherence, dietary indiscretion, myocardial ischemia, arrhythmias, infections, and anemia. Clinical features include dyspnea at rest, orthopnea, peripheral edema, and signs of pulmonary congestion. Diagnosis is based on clinical assessment, chest imaging, and natriuretic peptide levels. Management involves intravenous loop diuretics for decongestion, vasodilators for hypertensive crisis, and inotropic agents for cardiogenic shock; non-invasive positive pressure ventilation reduces the need for intubation. Prognosis depends on the underlying cause and response to therapy; early post-discharge follow-up and transition-of-care programs reduce readmission rates.

\subsection{Heart Failure with Preserved Ejection Fraction}
\label{sec:Heart Failure with Preserved Ejection Fraction}
Heart failure with preserved ejection fraction accounts for approximately half of heart failure cases and is characterized by symptoms of heart failure with left ventricular ejection fraction of 50\% or greater; it is associated with aging, hypertension, obesity, diabetes mellitus, and atrial fibrillation. The condition results from diastolic dysfunction (impaired myocardial relaxation and increased stiffness), elevated filling pressures, and systemic inflammation. Clinical features include exertional dyspnea, fatigue, and volume overload. Diagnosis is often challenging, requiring integration of symptoms, echocardiography (including diastolic function assessment), elevated natriuretic peptides, and sometimes invasive hemodynamic measurement. Management targets comorbidities, volume management with diuretics, and SGLT2 inhibitors (empagliflozin, dapagliflozin); exercise training and weight management are important non-pharmacological interventions. Prognosis is variable, with similar mortality rates to heart failure with reduced ejection fraction.

\subsection{Heart Failure with Reduced Ejection Fraction}
\label{sec:Heart Failure with Reduced Ejection Fraction}
Heart failure with reduced ejection fraction is a clinical syndrome characterized by symptoms and signs of heart failure with left ventricular ejection fraction of 40\% or less; common etiologies include ischemic heart disease, hypertension, dilated cardiomyopathy, valvular disease, and myocarditis. The condition results from neurohormonal activation (sympathetic nervous system, renin-angiotensin-aldosterone system) that drives disease progression. Clinical features range from dyspnea and fatigue (NYHA class II--IV) to volume overload with congestion. Diagnosis is based on clinical assessment and echocardiography demonstrating reduced left ventricular ejection fraction. Management involves four pillars of guideline-directed medical therapy: ACE inhibitor, ARB, or ARNI; beta-blocker; mineralocorticoid receptor antagonist; and SGLT2 inhibitor; device therapy with implantable cardioverter-defibrillator and cardiac resynchronization therapy reduces sudden death and improves symptoms. Prognosis has improved significantly with modern medical therapy.

%-------------------------------------------------------------
\section{Infective Endocarditis}
%-------------------------------------------------------------

%-------------------------------------------------------------

%-------------------------------------------------------------

%-------------------------------------------------------------

%-------------------------------------------------------------

%-------------------------------------------------------------

%-------------------------------------------------------------

\label{sec:section:Infective Endocarditis}
%-------------------------------------------------------------%-------------------------------------------------------------

\subsection{Infective Endocarditis}
\label{sec:Infective Endocarditis}
Infective endocarditis is an infection of the endocardial surface of the heart, most commonly the heart valves; acute infective endocarditis is caused by virulent organisms (\textit{Staphylococcus aureus}) and progresses rapidly, while subacute disease is caused by less virulent organisms (viridans group streptococci). Risk factors include prosthetic valves, intravenous drug use, structural heart disease, and indwelling catheters. The condition results from microbial infection of damaged or abnormal endocardium. Clinical features include fever, chills, night sweats, malaise, weight loss, new or changing cardiac murmur, splinter hemorrhages, Osler's nodes, Janeway lesions, Roth spots, and splenomegaly. Diagnosis requires modified Duke criteria: major criteria (positive blood cultures, evidence of endocardial involvement on echocardiography) and minor criteria (predisposing condition, fever, vascular phenomena, immunologic phenomena). Management involves prolonged intravenous antibiotics (4--6 weeks) based on culture and sensitivity; surgical intervention is indicated for heart failure, uncontrolled infection, large vegetations, or perivalvular abscess. Prognosis depends on the causative organism, patient factors, and timeliness of intervention.

%-------------------------------------------------------------
\section{Pericardial Disease}
%-------------------------------------------------------------

%-------------------------------------------------------------

%-------------------------------------------------------------

%-------------------------------------------------------------

%-------------------------------------------------------------

%-------------------------------------------------------------

%-------------------------------------------------------------

\label{sec:Pericardial Disease}
%-------------------------------------------------------------%-------------------------------------------------------------

\subsection{Cardiac Tamponade}
\label{sec:Cardiac Tamponade}
Cardiac tamponade is the accumulation of pericardial fluid under pressure, causing compression of the heart and impaired cardiac filling, leading to reduced cardiac output and obstructive shock; causes include malignancy (most common in developed countries), trauma, uremia, infection, aortic dissection, and post-procedural causes. The condition results from rapid or excessive pericardial fluid accumulation. Clinical features include Beck's triad (hypotension, muffled heart sounds, and jugular venous distension) and pulsus paradoxus. Diagnosis is by echocardiography, which shows pericardial effusion with diastolic collapse of the right atrium and right ventricle. Management is emergent pericardiocentesis or surgical pericardiotomy. Prognosis depends on the underlying cause and timeliness of intervention.

\subsection{Constrictive Pericarditis}
\label{sec:Constrictive Pericarditis}
Constrictive pericarditis is a condition in which the pericardium becomes thickened, fibrotic, and calcified, restricting diastolic filling of the heart; common causes include prior cardiac surgery, radiation therapy, tuberculosis, idiopathic pericarditis, and autoimmune disease. The condition results from chronic inflammation leading to pericardial fibrosis and restriction. The pathophysiology mimics restrictive cardiomyopathy but with ventricular interdependence and exaggerated respiratory variation in filling. Clinical features include signs of right heart failure (edema, ascites, hepatomegaly, elevated jugular venous pressure with Kussmaul's sign). Diagnosis involves echocardiography, cardiac MRI (pericardial thickening and late enhancement), and cardiac catheterization (equalization of diastolic pressures). Management is pericardiectomy, which is curative but carries significant surgical morbidity. Prognosis is excellent after successful pericardiectomy.

\subsection{Pericarditis}
\label{sec:Pericarditis}
Pericarditis is an inflammation of the pericardium, the fibrous sac surrounding the heart; acute pericarditis is the most common pericardial disease, with viral etiology being most frequent (Coxsackievirus, echovirus, adenovirus, SARS-CoV-2). Other causes include bacterial infection, uremia, autoimmune diseases, malignancy, and post-MI (Dressler syndrome). The condition results from inflammation of the pericardial layers. Clinical features include sharp, pleuritic chest pain that worsens with inspiration and improves with sitting forward, pericardial friction rub on auscultation, and diffuse ST-segment elevation on ECG. Diagnosis is based on clinical features, ECG, and echocardiography (may show pericardial effusion). Management includes NSAIDs (ibuprofen, aspirin) and colchicine for 3 months; corticosteroids are used for refractory or recurrent cases. Prognosis is generally good, with most cases resolving with treatment.

%-------------------------------------------------------------
\section{Peripheral Vascular Diseases}
%-------------------------------------------------------------

%-------------------------------------------------------------

%-------------------------------------------------------------

%-------------------------------------------------------------

%-------------------------------------------------------------

%-------------------------------------------------------------

%-------------------------------------------------------------

\label{sec:Peripheral Vascular Diseases}
%-------------------------------------------------------------%-------------------------------------------------------------

\subsection{Peripheral Arterial Disease}
\label{sec:Peripheral Arterial Disease}
Peripheral arterial disease is atherosclerotic narrowing of the arteries supplying the extremities, predominantly the lower limbs, resulting in reduced blood flow; it affects approximately 10--15\% of adults over 50 and is a marker of systemic atherosclerosis. The most important risk factor is cigarette smoking, followed by diabetes mellitus, hypertension, hyperlipidemia, and chronic kidney disease. The condition results from atherosclerotic plaque accumulation in peripheral arteries. Most patients are asymptomatic; intermittent claudication (cramping leg pain triggered by exertion and relieved by rest within 10 minutes) is the hallmark symptom; severe disease presents with critical limb ischemia (rest pain, non-healing ulcers, gangrene). Diagnosis is by the ankle-brachial index: a ratio of 0.90 or less is diagnostic. Management includes smoking cessation (most important intervention), antiplatelet therapy, statin therapy, supervised exercise programs, and management of comorbidities; revascularization is indicated for lifestyle-limiting claudication or critical limb ischemia. Prognosis depends on disease severity and risk factor modification.

%-------------------------------------------------------------
\section{Valvular Heart Disease}
%-------------------------------------------------------------

%-------------------------------------------------------------

%-------------------------------------------------------------

%-------------------------------------------------------------

%-------------------------------------------------------------

%-------------------------------------------------------------

%-------------------------------------------------------------

\label{sec:Valvular Heart Disease}
%-------------------------------------------------------------%-------------------------------------------------------------

\subsection{Aortic Dissection}
\label{sec:Aortic Dissection}
Aortic dissection is a life-threatening medical emergency caused by a tear in the aortic intima that allows blood to surge into the aortic media, creating a false lumen. The condition results from chronic hypertension, connective tissue disorders (Marfan syndrome, Ehlers-Danlos syndrome), or bicuspid aortic valve. Clinical features include sudden, severe, tearing or ripping chest pain radiating to the back (interscapular region), asymmetric blood pressure between arms, and neurological deficits. Diagnosis is confirmed urgently by CT angiography of the chest and abdomen or transesophageal echocardiography (TEE). Stanford Type A (ascending aorta) requires immediate emergency surgical repair; Stanford Type B (descending aorta) is managed medically with IV beta-blockers (esmolol, labetalol) to control heart rate and blood pressure unless complications occur. Prognosis is guarded, with high early mortality.

\subsection{Aortic Regurgitation}
\label{sec:Aortic Regurgitation}
Aortic regurgitation is incompetence of the aortic valve allowing diastolic backflow of blood from the aorta into the left ventricle; causes include bicuspid aortic valve, endocarditis, aortic root dilation (Marfan syndrome, aortic dissection, hypertension), and rheumatic disease. The condition results from aortic valve leaflet or aortic root abnormalities. The chronic volume overload leads to left ventricular dilation and eccentric hypertrophy, eventually causing heart failure. Patients may be asymptomatic for years despite severe regurgitation. Clinical signs include a wide pulse pressure, water-hammer pulse, decrescendo diastolic murmur, and de Musset's sign. Diagnosis is by echocardiography. Management involves surgical aortic valve repair or replacement, indicated for symptomatic patients or when left ventricular ejection fraction falls below 50\% or left ventricular end-systolic diameter exceeds 50 mm. Prognosis is excellent after surgical correction.

\subsection{Aortic Stenosis}
\label{sec:Aortic Stenosis}
Aortic stenosis is a narrowing of the aortic valve opening and the most common valvular disease in developed countries; degenerative calcific aortic stenosis is the most common etiology in older adults, while bicuspid aortic valve (present in 1--2\% of the population) is the most common cause in younger patients. The condition results from progressive calcification and restriction of the valve leaflets, causing pressure overload of the left ventricle, compensatory hypertrophy, and eventually heart failure. Clinical features include the classic triad of angina, syncope, and dyspnea. Diagnosis is by echocardiography showing severe aortic stenosis (aortic valve area less than 1.0 cm$^2$, mean gradient greater than 40 mmHg). Management is aortic valve replacement with surgical (SAVR) or transcatheter (TAVR) approaches; TAVR is increasingly preferred in intermediate- and high-risk patients. Prognosis is poor without intervention (mortality 50\% at 2 years after symptom onset) but excellent after successful valve replacement.

\subsection{Coarctation of the Aorta}
\label{sec:Coarctation of the Aorta}
Coarctation of the Aorta is a congenital cardiac anomaly characterized by a discrete constriction of the aortic lumen, most commonly located distal to the origin of the left subclavian artery near the insertion of the ductus arteriosus (juxtaductal). The condition results from ectopic ductal tissue tissue contraction or impaired embryonic aortic arch development. Clinical features include upper extremity hypertension with delayed, weak, or absent femoral pulses and a lower extremity blood pressure gradient ($>10--20$ mmHg lower in legs); patients may present with headache, cold feet, leg claudication, or heart failure. A continuous systolic murmur is best heard over the back. Chest radiograph demonstrates rib notching (due to collateral intercostal arterial enlargement) and the 'figure-3' sign on the aorta. Diagnosis is confirmed by echocardiography or cardiac MRI. Management involves surgical resection with end-to-end anastomosis or balloon angioplasty with stenting. Prognosis is good following early intervention.

\subsection{Hypertension}
\label{sec:Hypertension}
Hypertension is defined as a sustained elevation of systemic arterial blood pressure (systolic $\ge$130 mmHg or diastolic $\ge$80 mmHg per ACC/AHA guidelines). Essential (primary) hypertension accounts for 90--95\% of cases and results from complex interactions between genetic factors, renal sodium retention, sympathetic nervous system hyperreactivity, and vascular smooth muscle hypertrophy. Secondary hypertension (5--10\% of cases) results from renal parenchymal disease, renovascular hypertension, primary aldosteronism, obstructive sleep apnea, or endocrine tumors. Most patients are asymptomatic until end-organ damage occurs. Diagnosis requires average blood pressure measurements from $\ge$2 occasions. Management includes lifestyle modification (DASH diet, sodium restriction, exercise) and first-line pharmacotherapy (ACE inhibitors/ARBs, thiazide diuretics, calcium channel blockers). Prognosis is excellent with early control.

\subsection{Mitral Regurgitation}
\label{sec:Mitral Regurgitation}
Mitral regurgitation (MR) is the retrograde flow of blood from the left ventricle into the left atrium during systole. Primary (organic) MR results from intrinsic valve apparatus pathology, including mitral valve prolapse (myxomatous degeneration), infective endocarditis, or rheumatic heart disease; secondary (functional) MR results from left ventricular remodeling or dilation (ischemic or dilated cardiomyopathy) displacing papillary muscles. Clinical features include exertional dyspnea, fatigue, orthopnea, and atrial fibrillation; physical examination reveals a high-pitched holosystolic (pansystolic) murmur at the apex radiating to the left axilla, S3 gallop, and displaced hyperdynamic apical impulse. Diagnosis is confirmed by transthoracic echocardiography assessing valve anatomy, regurgitant jet volume, and left ventricular function. Management of severe MR includes medical optimization (afterload reduction with ACE inhibitors/ARBs, diuretics) and timely surgical mitral valve repair or replacement prior to irreversible left ventricular dysfunction. Prognosis depends on baseline LV ejection fraction.

\subsection{Mitral Stenosis}
\label{sec:Mitral Stenosis}
Mitral stenosis is a narrowing of the mitral valve orifice, most commonly caused by rheumatic heart disease, with a female predominance. The condition results from progressive fibrosis and fusion of the mitral commissures, leading to elevated left atrial pressure, pulmonary congestion, and eventually pulmonary hypertension and right heart failure. The normal mitral valve area is 4--6 cm$^2$; symptoms typically appear when the area narrows to less than 2 cm$^2$. Clinical features include dyspnea, orthopnea, hemoptysis, and atrial fibrillation. Diagnosis is by echocardiography. Management includes percutaneous mitral balloon valvulotomy for suitable anatomy, surgical commissurotomy, or mitral valve replacement; anticoagulation is indicated for atrial fibrillation or thromboembolism. Prognosis depends on the severity and success of intervention.

\subsection{Mitral Valve Prolapse}
\label{sec:Mitral Valve Prolapse}
Mitral valve prolapse is the most common valvular abnormality, affecting 2--3\% of the population, and involves billowing of one or both mitral valve leaflets into the left atrium during systole, with or without mitral regurgitation. Most cases are idiopathic or associated with myxomatous degeneration; mitral valve prolapse is associated with Marfan syndrome and Ehlers-Danlos syndrome. The condition results from structural abnormality of the mitral valve apparatus. Most patients are asymptomatic and discovered incidentally; clinical features may include palpitations, chest pain, and dysautonomic symptoms. Diagnosis is by echocardiography. Asymptomatic patients with minimal regurgitation require only reassurance; severe regurgitation may require surgical repair. Prognosis is excellent in most cases; complications include mitral regurgitation, infective endocarditis, and rarely, arrhythmias.

\subsection{Peripheral Artery Disease}
\label{sec:Peripheral Artery Disease}
Peripheral artery disease (PAD) is a chronic atherosclerotic occlusion of systemic arteries supplying the lower extremities. The condition results from endothelial dysfunction and atheroma accumulation reducing tissue perfusion. Risk factors include cigarette smoking, diabetes mellitus, hypertension, hyperlipidemia, and advanced age. Clinical features include intermittent claudication (reproducible lower extremity muscle pain precipitated by exertion and relieved by rest), cold extremities, hair loss, thickened toenails, and resting pain or ischemic ulceration/gangrene (critical limb ischemia). Diagnosis is confirmed by an Ankle-Brachial Index (ABI) $<0.90$, with duplex ultrasonography or CT angiography defining anatomical stenosis. Management includes lifestyle modification (smoking cessation, structured exercise program), antiplatelet therapy (aspirin or clopidogrel), statin therapy, cilostazol for claudication, and endovascular angioplasty/stenting or surgical bypass for severe or refractory disease. Prognosis is strongly correlated with coexisting coronary and cerebrovascular disease.

\subsection{Rheumatic Heart Disease}
\label{sec:Rheumatic Heart Disease}
Rheumatic heart disease (RHD) is the chronic valvular damage resulting from acute rheumatic fever, an autoimmune response to group A streptococcal pharyngitis. The condition results from molecular mimicry between streptococcal M protein and human cardiac tissue, leading to pancarditis and chronic fibrotic scarring of heart valves; the mitral valve is affected in over 90\% of cases, leading to mitral stenosis or regurgitation. Clinical features include exertional dyspnea, fatigue, palpitations, and signs of heart failure or stroke secondary to atrial fibrillation. Diagnosis is based on echocardiography showing valvular thickening, commissural fusion, and regurgitation/stenosis. Management includes secondary antibiotic prophylaxis (IM benzathine penicillin G every 3--4 weeks) to prevent recurrent streptococcal infections, medical management of heart failure, and surgical valve repair or replacement. Prognosis depends on valve damage severity.

\subsection{Tricuspid Regurgitation}
\label{sec:Tricuspid Regurgitation}
Tricuspid regurgitation is backflow of blood through the tricuspid valve during systole; functional tricuspid regurgitation (secondary to right ventricular dilation from pulmonary hypertension, left-sided heart disease, or atrial fibrillation) is the most common type. Primary tricuspid regurgitation results from endocarditis, rheumatic disease, carcinoid syndrome, or pacemaker lead infection. The condition results from structural or functional abnormalities of the tricuspid valve apparatus. Clinical features of severe tricuspid regurgitation include right heart failure with peripheral edema, ascites, hepatomegaly, and pulsatile jugular venous distension. Diagnosis is by echocardiography with Doppler. Management targets the underlying cause; surgical valve repair or replacement is considered for severe symptomatic disease. Prognosis depends on the underlying cause and severity.

